% Problem-Led Reading
% An open, modular reading template with an optional AI margin.
% Compile with XeLaTeX or LuaLaTeX.
\documentclass[10pt,fontset=none]{ctexart}

\usepackage{iftex}
\ifPDFTeX\errmessage{Please compile with XeLaTeX or LuaLaTeX}\fi
\usepackage{fontspec}
\usepackage{geometry}
\usepackage{xcolor}
\usepackage{microtype}
\usepackage{hyperref}
\usepackage{fancyhdr}
\usepackage[most]{tcolorbox}
\usepackage{amsmath,amssymb,mathtools}
\usepackage{tabularx,array}

\geometry{
  paperwidth=7.5in,paperheight=10in,
  left=0.62in,right=0.62in,top=0.48in,bottom=0.45in,
  headsep=0in,footskip=0.23in
}

\IfFontExistsTF{Songti SC}
  {\setCJKmainfont{Songti SC}[AutoFakeBold=2.0]}
  {\IfFontExistsTF{Source Han Serif SC}
    {\setCJKmainfont{Source Han Serif SC}[AutoFakeBold=2.0]}
    {\setCJKmainfont{FandolSong-Regular.otf}[BoldFont=FandolSong-Bold.otf]}}
\IfFontExistsTF{Times New Roman}
  {\setmainfont{Times New Roman}}
  {\setmainfont{TeX Gyre Termes}}
\IfFontExistsTF{Helvetica Neue}
  {\setsansfont{Helvetica Neue}}
  {\setsansfont{TeX Gyre Heros}}
\IfFontExistsTF{Hiragino Sans GB}
  {\setCJKsansfont{Hiragino Sans GB}[AutoFakeBold=2.0]}
  {\setCJKsansfont{FandolHei-Regular.otf}[BoldFont=FandolHei-Bold.otf]}

\definecolor{navy}{HTML}{163A60}
\definecolor{ink}{HTML}{171B20}
\definecolor{muted}{HTML}{707981}
\definecolor{rule}{HTML}{D3D8DC}
\definecolor{soft}{HTML}{F4F6F7}
\definecolor{paper}{HTML}{FCFBF8}

\pagecolor{paper}
\color{ink}
\linespread{1.045}
\setlength{\parindent}{0pt}
\setlength{\parskip}{0pt}
\setlength{\abovedisplayskip}{4pt}
\setlength{\belowdisplayskip}{4pt}
\raggedbottom
\clubpenalty=10000
\widowpenalty=10000

% -----------------------------------------------------------------
% Reuse: change only the example source identity.
% -----------------------------------------------------------------
\newcommand{\SeriesNumber}{01}
\newcommand{\SourceTitle}{Probability Inequalities for Sums of Bounded Random Variables}
\newcommand{\SourceAuthors}{Wassily Hoeffding}
\newcommand{\SourceContext}{JASA 58 (1963)}
\newcommand{\SourceURL}{https://doi.org/10.1080/01621459.1963.10500830}
\newcommand{\ReadDate}{2026.08.18}

\hypersetup{
  colorlinks=true,urlcolor=navy,linkcolor=navy,
  pdfauthor={MIStatlE},
  pdftitle={Problem-Led Reading}
}

\fancyhf{}
\renewcommand{\headrulewidth}{0pt}
\fancyfoot[L]{\footnotesize\sffamily\color{muted}PROBLEM-LED READING \enspace \SeriesNumber}
\fancyfoot[R]{\footnotesize\sffamily\color{navy}@MIStatlE \enspace\textperiodcentered\enspace \thepage/4}
\pagestyle{fancy}

% -----------------------------------------------------------------
% A small grammar of freely composable note elements.
% -----------------------------------------------------------------
\newcommand{\PageMark}[2]{%
  \noindent
  \begin{minipage}[t]{0.72\linewidth}
    {\sffamily\bfseries\scriptsize\color{navy}#1}
  \end{minipage}\hfill
  \begin{minipage}[t]{0.24\linewidth}
    \raggedleft{\sffamily\scriptsize\color{muted}#2}
  \end{minipage}}

\newcommand{\PageTitle}[2]{%
  \PageMark{PROBLEM-LED READING \enspace/\enspace #1}{\ReadDate}
  \vspace{0.18em}
  {\bfseries\fontsize{23pt}{27pt}\selectfont\color{ink}#2\par}}

\newcommand{\SourceLine}{%
  {\small\color{muted}\SourceTitle\par}
  \vspace{0.07em}
  {\footnotesize\color{muted}\SourceAuthors\quad\textperiodcentered\quad
   \SourceContext\quad\textperiodcentered\quad\href{\SourceURL}{source}\par}}

\newcommand{\RuleHead}[2]{%
  \vspace{0.56em}
  \noindent{\sffamily\bfseries\scriptsize\color{muted}#1}\hspace{0.65em}%
  {\sffamily\bfseries\large\color{navy}#2}\hspace{0.75em}%
  {\color{rule}\leaders\hrule height 0.45pt\hfill\kern0pt}\par
  \vspace{0.28em}}

\newcommand{\Tag}[1]{%
  {\sffamily\bfseries\scriptsize\color{navy}\MakeUppercase{#1}}}

\newcommand{\SourceNote}[3]{%
  \noindent
  \begin{minipage}[t]{0.17\linewidth}
    \vspace{0pt}{\sffamily\scriptsize\color{muted}#1}
  \end{minipage}
  \begin{minipage}[t]{0.17\linewidth}
    \vspace{0pt}\Tag{#2}
  \end{minipage}
  \begin{minipage}[t]{0.62\linewidth}
    \vspace{0pt}\raggedright #3
  \end{minipage}\par}

\newcommand{\MarginQuestion}[2]{%
  \Tag{#1}\par
  \vspace{0.12em}{\small #2}\par}

\newcommand{\OpenLine}[1]{%
  {\small\color{muted}#1}\par
  \vspace{0.32em}{\color{rule}\hrule height 0.45pt}\vspace{0.52em}}

\tcbset{
  openbox/.style={enhanced,colback=paper,colframe=rule,boxrule=0.45pt,
    arc=0pt,left=10pt,right=10pt,top=7.5pt,bottom=7.5pt},
  marginbox/.style={enhanced,colback=soft,frame hidden,
    borderline west={1.8pt}{0pt}{navy},arc=0pt,
    left=10pt,right=9pt,top=7.5pt,bottom=7.5pt}
}

\begin{document}

% =================================================================
% PAGE 1 -- blank reusable page
% =================================================================
\PageMark{PROBLEM-LED READING \enspace/\enspace \SeriesNumber}{OPEN TEMPLATE}

\vspace{0.40em}
{\bfseries\fontsize{28pt}{32pt}\selectfont\color{ink}问题研读\par}
\vspace{0.05em}
{\large\color{navy}Problem-Led Reading\par}

\vspace{0.48em}
{\small\color{muted}
  先读原文，形成初步判断；再让 AI 展开候选，最后回到来源修订。}

\RuleHead{START}{Reading question}

\begin{tcolorbox}[openbox]
  {\color{muted}阅读这份材料，我希望厘清：}\par
  \vspace{0.55em}{\color{rule}\hrule height 0.45pt}\vspace{0.62em}
  {\color{rule}\hrule height 0.45pt}
\end{tcolorbox}

\vspace{0.30em}
\begin{minipage}[t]{0.69\linewidth}
  \vspace{0pt}
  \RuleHead{SOURCE}{First pass}
  \OpenLine{\textsf{p./sec./fig.}\quad 一条可定位的主张、证据或机制。}
  \OpenLine{\textsf{?}\quad 一个阻碍理解或可能改变判断的问题。}

  \RuleHead{MY READ}{Current interpretation}
  \begin{tcolorbox}[openbox,height=1.10in,valign=top]
    {\small\color{muted}我目前如何理解它？哪一处仍不确定？}
  \end{tcolorbox}
\end{minipage}\hfill
\begin{minipage}[t]{0.26\linewidth}
  \vspace{0pt}
  \RuleHead{OPTIONAL}{AI margin}
  \begin{tcolorbox}[marginbox]
    \MarginQuestion{clarify}{隐含前提是什么？}
    \vspace{0.62em}
    \MarginQuestion{alternate}{替代解释是什么？}
    \vspace{0.62em}
    \MarginQuestion{check}{如何回到来源核验？}
  \end{tcolorbox}
\end{minipage}

\RuleHead{WORK}{Think on paper}

\begin{tcolorbox}[openbox,height=2.25in,valign=top]
  {\footnotesize\color{muted}
    quote / diagram / derivation / mechanism / objection}
\end{tcolorbox}

\RuleHead{RETURN}{Next place to resume}
\OpenLine{下次从以下位置继续：}

\newpage

% =================================================================
% PAGE 2 -- a raw first read before AI
% =================================================================
\PageTitle{ESSAY EXAMPLE / 01 FIRST PASS}{初读：形成判断}
\vspace{0.08em}
\SourceLine

\RuleHead{QUESTION}{Initial question}

{\bfseries\fontsize{17pt}{22pt}\selectfont
  为什么有界独立变量之和会呈现二次指数衰减？\par}

\vspace{0.16em}
{\small\color{muted}
  先处理原文，区分主张、证据与个人解释；暂不引入 AI 生成内容。}

\vspace{0.40em}
\begin{minipage}[t]{0.69\linewidth}
  \vspace{0pt}
  \RuleHead{TRACE}{Source notes}

  \SourceNote{定理}{claim}{
    若独立变量 $X_i\in[a_i,b_i]$，则对 $t>0$，
    \[
      \Pr\!\left(S_n-\mathbb ES_n\ge t\right)
      \le \exp\!\left(-\frac{2t^2}{\sum_i(b_i-a_i)^2}\right).
    \]}
  \vspace{0.55em}{\color{rule}\hrule height 0.4pt}\vspace{0.52em}
  \SourceNote{引理}{local bound}{
    对单个有界变量，
    $\log\mathbb E e^{\lambda(X_i-\mathbb EX_i)}
    \le \lambda^2(b_i-a_i)^2/8$。}
  \vspace{0.55em}{\color{rule}\hrule height 0.4pt}\vspace{0.52em}
  \SourceNote{证明}{mechanism}{
    独立性将和的 MGF 分解为乘积；再用 Markov--Chernoff 界并优化 $\lambda$。}

  \RuleHead{MY READ}{A tension worth keeping}
  \begin{tcolorbox}[openbox]
    {\bfseries 高斯型尾部似乎主要来自变量有界。}\par
    \vspace{0.35em}
    {\small\color{muted}
      个人解释；尚未区分有界性与独立性的作用。}
  \end{tcolorbox}
\end{minipage}\hfill
\begin{minipage}[t]{0.26\linewidth}
  \vspace{0pt}
  \RuleHead{AI OFF}{First pass}
  \begin{tcolorbox}[marginbox]
    {\footnotesize\color{muted}
      暂不生成结论，仅保留初读中的张力与问题。}

    \vspace{0.72em}
    \MarginQuestion{concept}{二次指数从哪一步出现？}
    \vspace{0.72em}
    \MarginQuestion{evidence}{独立性只方便证明，还是决定尺度？}
    \vspace{0.72em}
    \MarginQuestion{boundary}{变量完全相关时会怎样？}
  \end{tcolorbox}
\end{minipage}

\vfill
\RuleHead{STOP}{First-pass record}
{\small
  一个定理、一个引理、一条证明链、三个问题。初读到此为止。}
\vspace{0.24in}

\newpage

% =================================================================
% PAGE 3 -- AI expands the search space without owning the note
% =================================================================
\PageTitle{ESSAY EXAMPLE / 02 EXPAND}{扩展：生成候选}
\vspace{0.08em}
\SourceLine

\begin{minipage}[t]{0.69\linewidth}
  \vspace{0pt}
  \RuleHead{HANDOFF}{Inputs and constraints}
  \begin{tcolorbox}[openbox]
    {\small
      输入：定理、MGF 引理、证明链与初步判断。\par
      \color{muted}约束：只拆分假设的作用；反例与推广必须明确标记。}
  \end{tcolorbox}

  \RuleHead{BRANCH 01}{Separate the roles}
  {\bfseries 候选解释：}\;
  有界性给出单变量的二次 MGF 控制；独立性使对数 MGF 可加。
  两者共同产生平方和尺度。\par
  \vspace{0.25em}
  {\footnotesize\color{muted}
    核验：逐步标出每个假设进入证明的位置。}

  \RuleHead{BRANCH 02}{Break independence}
  {\bfseries 候选反例：}\;
  令 $Y$ 为 Rademacher 变量且 $X_i=Y$。各变量均有界，但 $S_n=nY$；
  在 $t=n$ 处真实概率为 $1/2$，独立情形的指数界失效。\par
  \vspace{0.25em}
  {\footnotesize\color{muted}
    结论：独立性不是可以无条件删除的技术细节。}

  \RuleHead{BRANCH 03}{Transfer carefully}
  {\bfseries 候选迁移：}\;
  对鞅差分，以条件 MGF 控制替代乘积分解，得到 Azuma--Hoeffding 型界。\par
  \vspace{0.25em}
  {\footnotesize\color{muted}
    标记：这是后续推广，不属于原定理本身。}
\end{minipage}\hfill
\begin{minipage}[t]{0.26\linewidth}
  \vspace{0pt}
  \RuleHead{OPTIONAL}{AI margin}
  \begin{tcolorbox}[marginbox]
    {\footnotesize\color{muted}
      AI 仅扩展假设空间，不修改主笔记。}

    \vspace{0.70em}
    \MarginQuestion{clarify}{分别指出有界性与独立性的作用。}
    \vspace{0.70em}
    \MarginQuestion{alternate}{方差控制会导向什么替代界？}
    \vspace{0.70em}
    \MarginQuestion{falsify}{构造删除独立性后的反例。}
    \vspace{0.70em}
    \MarginQuestion{provenance}{区分原定理与鞅推广。}
  \end{tcolorbox}
\end{minipage}

\vfill
\RuleHead{SELECT}{Selected branch}
\begin{tcolorbox}[openbox]
  推进分支 02：用完全相关的 Rademacher 变量检验独立性是否只是技术条件。
  推广分支暂不写回。
\end{tcolorbox}
\vspace{0.22in}

\newpage

% =================================================================
% PAGE 4 -- write back only what survived checking
% =================================================================
\PageTitle{ESSAY EXAMPLE / 03 CHECK}{回源：修订判断}
\vspace{0.30em}
\SourceLine

\RuleHead{KEEP}{Source-backed}
\begin{tcolorbox}[openbox,fontupper=\small]
  原定理同时使用每个变量的取值范围与独立性；
  尾部尺度由 $\sum_i(b_i-a_i)^2$ 控制。
\end{tcolorbox}

\RuleHead{REVISE}{Before / after}
\begin{minipage}[t]{0.48\linewidth}
  \vspace{0pt}\Tag{before}\par
  \vspace{0.25em}
  \begin{tcolorbox}[openbox,height=1.05in,valign=top,fontupper=\small]
    高斯型尾部主要来自变量有界。
  \end{tcolorbox}
\end{minipage}\hfill
\begin{minipage}[t]{0.48\linewidth}
  \vspace{0pt}\Tag{after}\par
  \vspace{0.25em}
  \begin{tcolorbox}[openbox,height=1.05in,valign=top,fontupper=\small]
    有界性提供局部 MGF 控制；独立性使控制量可加。二者缺一不可。
  \end{tcolorbox}
\end{minipage}

\RuleHead{OPEN}{One question to carry}
\begin{tcolorbox}[openbox,fontupper=\small]
  独立性可以在多大程度上被条件 MGF 控制、鞅差分或弱依赖结构替代？
  这是从原证明结构延伸出的新问题。
\end{tcolorbox}

\vfill
\RuleHead{AI ROLE}{Expand, then step back}
\begin{tcolorbox}[marginbox]
  {\small
    AI 提供候选解释；最终写回只依据原文核验与个人判断修订。}
\end{tcolorbox}
\vspace{0.24in}

\end{document}
